% Paper 1, §12 — Inference adjuncts
% Anchors: kb/notes/inference/{kv-cache-management,quantization,structured-output}.md
%          kb/excerpts/{kwon2023,frantar2022,xiao2023-smoothquant,
%                       ma2024-bitnet,willard2023-outlines,xgrammar2024}.md

\section{Inference adjuncts}
\label{sec:inference-adjuncts}

The architectural choices catalogued in \cref{sec:tokenization-embedding}
through \cref{sec:multimodal} fix the model's mathematical surface; they
do not by themselves determine what the deployed system can serve at
matched cost. Three categories of \emph{inference-time} concern --- the
KV-cache lifecycle and its memory-allocation discipline, weight /
activation / KV quantization regimes, and structured-output enforcement
through \glsterm{constrained-decoding-2}{constrained decoding} --- are the load-bearing reason a 2026
frontier model is feasible to ship. Each constrains which choices in
\cref{sec:kv-sharing,sec:hardware-attention,sec:ffn,sec:state-space}
can be taken: paged allocation requires a block-friendly attention
kernel, ternary pretraining requires LLaMA-alike normalization and
gating, FSM-masked decoding requires an LM-head whose logit space is
tokenizer-aligned. The thesis of this section is that these are
\emph{architectural concerns, not bolt-ons}.

\subsection{KV-cache lifecycle and PagedAttention}
\label{sec:kv-paged}

The \glsterm{autoregressive}{autoregressive} decode loop reads, at every step, the cached K/V
projection of every prior token at every layer. For an OPT-13B request
of 2048 tokens the \glsterm{kv-cache}{KV cache} is $\approx 1.6$~GB, with the per-token
cost $2 \cdot L \cdot n_h \cdot d_h \cdot b$ bytes where $b$ is the
element width \citep[\S 3]{kwon2023}.\footnote{KB excerpt:
\texttt{kb/excerpts/kwon2023.md\#sec-3}.} Pre-2023 serving stacks
(FasterTransformer, Orca) pre-allocated each request a contiguous slab
sized to its \emph{maximum} sequence length, which produced three
sources of waste \citep[\S 3]{kwon2023}: \emph{reservation} (slots
held for future tokens not yet generated), \emph{internal
fragmentation} (slack between actual and maximum length, realized only
at request completion), and \emph{external fragmentation} (allocator
holes between slabs). Empirically the contiguous regime put only
$20.4$--$38.2$\% of allocated KV memory onto actual token state
\citep[\S 3, Fig.~2]{kwon2023}.

\citet{kwon2023} resolves the waste taxonomy by partitioning each
sequence's \glsterm{kv-cache}{KV cache} into fixed-size \emph{KV blocks} of $B$ tokens
(typical $B = 16$) and decomposing the attention sum block-wise
\citep[\S 4.1, Eq.~4]{kwon2023}\footnote{KB excerpt:
\texttt{kb/excerpts/kwon2023.md\#sec-4-1}.}:
\begin{equation}
A_{ij} = \frac{\exp(q_i^\top K_j / \sqrt{d})}{\sum_{t=1}^{\lceil i/B \rceil} \exp(q_i^\top K_t \mathbf{1} / \sqrt{d})}, \qquad
o_i = \sum_{j=1}^{\lceil i/B \rceil} V_j A_{ij}^\top,
\label{eq:pagedattn}
\end{equation}
with $K_j = (k_{(j-1)B+1}, \dots, k_{jB})$, $V_j = (v_{(j-1)B+1}, \dots, v_{jB})$,
and the kernel reading each block from a possibly non-contiguous
physical address. The arithmetic is the same as the canonical
scaled-dot-product attention of \cref{sec:hardware-attention}; the
contribution is layout, not arithmetic.

A \emph{block table} per request maps logical KV blocks to physical KV
blocks, exactly the role of an OS page table
\citep[\S 4.2]{kwon2023}.\footnote{KB excerpt:
\texttt{kb/excerpts/kwon2023.md\#sec-4-2}.} Three OS idioms transfer
unchanged: allocate-on-demand (no maximum-length pre-reservation),
copy-on-write at the block boundary for parallel sampling and beam
search \citep[\S 4.4]{kwon2023}, and swap-or-recompute for preemption
\citep[\S 4.5]{kwon2023}. The reported headline is $2$--$4\times$
throughput at matched latency vs.\ FasterTransformer/Orca on
OPT-13B/66B/175B \citep[abstract, \S 6.1]{kwon2023}; the cache-utilization
number rises from $\sim 30$\% to $96.3$\% \citep[\S 3, Fig.~2]{kwon2023}.

\begin{analogy}
\glsterm{pagedattention}{PagedAttention} is virtual memory for the \glsterm{kv-cache}{KV cache}. The block table is
a page table; allocate-on-demand is demand paging; copy-on-write across
sampled continuations is the same pattern as \texttt{fork()}'s
COW pages; swap evicts blocks to CPU RAM as a swap device. Returning
to math: the block decomposition in \cref{eq:pagedattn} is the
canonical attention sum re-indexed over physical blocks rather than
contiguous positions; the kernel's flexibility to fetch each $K_j, V_j$
from a non-contiguous address is what frees the allocator from the
maximum-length reservation that produced the $\sim 70$\% waste.
\end{analogy}

The architectural pressure runs the other way too: \glsterm{pagedattention}{PagedAttention}'s
block-decomposed kernel is what enables the cross-request prefix sharing
of \glsterm{radixattention}{RadixAttention} \citep[\S 3]{zheng2024-sglang}\footnote{KB excerpt:
\texttt{kb/excerpts/zheng2024-sglang.md\#sec-3}.}, the
training-time KV compression of \glsterm{mqa}{MQA} / \glsterm{gqa}{GQA} / \glsterm{multi-head-latent-attention}{MLA} in
\cref{sec:kv-sharing} (which act on the per-token block size), and the
KV-quantization regimes below (which change the element width $b$ but
not the layout).

\subsection{Quantization regimes}
\label{sec:quantization}

\glsterm{llm}{LLM} quantization is a family of decisions indexed by \emph{what} is
quantized (weights, activations, \glsterm{kv-cache}{KV cache}), \emph{when} (post-training
vs.\ training-aware vs.\ native low-bit pretraining), and \emph{to what
precision} (8-bit, 4-bit, 2-bit, ternary). The uniform integer map
underlying every method below is
\begin{equation}
\bar{X}^{\text{INT}N} = \left\lceil \frac{X^{\text{FP16}}}{\Delta} \right\rfloor,
\qquad
\Delta = \frac{\max(|X|)}{2^{N-1} - 1},
\label{eq:uniform-quant}
\end{equation}
with $N$-bit signed levels and step size $\Delta$
\citep[\S 2, Eq.~1]{xiao2023-smoothquant}.\footnote{KB excerpt:
\texttt{kb/excerpts/xiao2023-smoothquant.md\#sec-2}.} Hardware
constrains the granularity: tensor-core GEMMs apply scale factors only
along outer dimensions (token $T$ or output channel $C_o$), never along
the inner contraction dimension $C_i$
\citep[\S 3]{xiao2023-smoothquant}. Per-channel-of-activations is
infeasible without re-engineering the kernel. This single constraint
generates the lineage of methods.

\paragraph{INT8 (W8A8) --- SmoothQuant.}
Past $\sim 6.7$B \glsterm{parameters}{parameters}, \glsterm{llm}{LLM} activations develop systematic
outliers $\sim 100\times$ larger than typical values, persistent in
fixed channels across tokens \citep[\S 3]{xiao2023-smoothquant}. Naive
per-tensor INT8 catastrophically degrades quality (OPT-175B accuracy
$71.6\%$ FP16 $\to$ $32.3\%$ per-tensor INT8
\citep[\S 5.2, Table~1]{xiao2023-smoothquant}).\footnote{KB excerpt:
\texttt{kb/excerpts/xiao2023-smoothquant.md\#sec-5-2}.}
\citet{xiao2023-smoothquant} resolves this with a \emph{migration}
trick that is mathematically a free rescaling
\citep[\S 4, Eq.~3]{xiao2023-smoothquant}:
\begin{equation}
Y = (X \operatorname{diag}(s)^{-1}) \cdot (\operatorname{diag}(s) W) = \hat{X}\hat{W},
\label{eq:smoothquant-eq}
\end{equation}
with the per-channel smoothing factor $s_j$ chosen to balance the
quantization difficulty between activations and weights via a
migration-strength hyperparameter $\alpha$
\citep[\S 4, Eq.~4]{xiao2023-smoothquant}:
\begin{equation}
s_j = \frac{\max(|X_j|)^{\alpha}}{\max(|W_j|)^{1-\alpha}},
\qquad \alpha = 0.5 \text{ default}.
\label{eq:smoothquant-s}
\end{equation}
At $\alpha = 0$ the difficulty stays on activations (broken); at
$\alpha = 1$ it migrates entirely to weights (also broken because
weights now have the outlier shape they previously lacked); $\alpha
\approx 0.5$ is the sweet spot. Crucially, $s$ is fused offline into
the preceding linear or normalization layer so the runtime cost is
zero, and per-tensor INT8 then recovers the FP16 baseline ($71.4\%$ on
OPT-175B \citep[\S 5.2, Table~1]{xiao2023-smoothquant}).

\paragraph{INT4 (W4A16) --- GPTQ.}
For weight-only INT4, \citet{frantar2022} adapt the layer-wise
reconstruction objective \citep[\S 3, Eq.~1]{frantar2022}\footnote{KB
excerpt: \texttt{kb/excerpts/frantar2022.md\#sec-3-layerwise}.}
\begin{equation}
\operatorname*{argmin}_{\widehat{W}} \, \|W X - \widehat{W} X\|_2^2
\label{eq:gptq-objective}
\end{equation}
to \glsterm{llm}{LLM} scale by accelerating the Optimal Brain Quantization baseline
along three axes: arbitrary fixed column order rather than per-row
greedy ordering, lazy batched updates over $B = 128$ columns, and a
Cholesky reformulation of the inverse-Hessian update with damping
$\lambda \approx 1\%$ of the diagonal \citep[\S 4]{frantar2022}.\footnote{KB
excerpt: \texttt{kb/excerpts/frantar2022.md\#sec-4-algorithm}.} The
joint effect: OPT-175B and BLOOM-176B at 3--4 bits/weight track FP16
\glsterm{perplexity}{perplexity}, where round-to-nearest at the same precision diverges
($110$ and $571$ respectively
\citep[\S 1, Fig.~1]{frantar2022}).\footnote{KB excerpt:
\texttt{kb/excerpts/frantar2022.md\#sec-1-figure-1}.} \glsterm{gptq}{GPTQ} is the
canonical accurate one-shot weight-only PTQ method at the 100B+ scale
and is the reference against which \glsterm{awq}{AWQ} and the GGUF Q$k$\_$\!$M family
are still benchmarked.

\paragraph{FP8 (training).}
Training in \glsterm{fp8}{FP8} with tile-based scales (128$\times$128 weight blocks,
1$\times$128 activation vectors) was first validated at trillion-token
scale by \glsterm{precision-pinning}{DeepSeek-V3} \citep[\S 2.5]{deepseek-v3}, with H100/H200/Blackwell
hardware \glsterm{fp8}{FP8} paths. We treat \glsterm{fp8}{FP8} in
\cref{sec:training-mixed-precision-paper-2}; the relevance to this
section is that an \glsterm{fp8}{FP8}-trained checkpoint can be served \emph{without}
PTQ, collapsing two stages into one.

\paragraph{NVFP4 (frontier 2026).}
Microscaling formats from the OCP standard pair small blocks with a
shared exponent. \glsterm{nvfp4}{NVFP4} is NVIDIA's Blackwell-generation 4-bit
floating-point with per-microblock scales and is the first 4-bit format
shipped with native tensor-core acceleration; vendor materials report
it as production-ready for inference at frontier scale, though
third-party reproductions outside NVIDIA-published benchmarks are
sparse as of 2026-Q2 \citep[FORUM-SIGNAL]{deepseek-v3}.\footnote{No
peer-reviewed primary source for \glsterm{nvfp4}{NVFP4} \glsterm{llm}{LLM} quality at frontier scale
exists in the KB as of \texttt{2026-05}; the citation here points to
the closest analogue (\glsterm{fp8}{FP8} production training) and flags the gap.}

\paragraph{1.58-bit --- BitNet.}
\citet{ma2024-bitnet} take the dual position: rather than quantize a
trained model, train from scratch with ternary weights. Every weight
is constrained to $\{-1, 0, +1\}$ ($\log_2 3 \approx 1.585$~bits) via
an absmean RoundClip \citep[\S 2, Eq.~1--3]{ma2024-bitnet}\footnote{KB
excerpt: \texttt{kb/excerpts/ma2024-bitnet.md\#sec-2}.}:
\begin{equation}
\widetilde{W} = \operatorname{RoundClip}\!\left(\frac{W}{\gamma + \epsilon}, -1, 1\right),
\qquad
\gamma = \frac{1}{nm} \sum_{ij} |W_{ij}|,
\label{eq:bitnet-quant}
\end{equation}
with $\operatorname{RoundClip}(x, a, b) = \max(a, \min(b, \operatorname{round}(x)))$;
activations remain INT8 (the regime is W1.58A8, not pure 1-bit). The
architecture is otherwise LLaMA-alike (\glsterm{rmsnorm}{RMSNorm}, \glsterm{swiglu}{SwiGLU}, \glsterm{rope}{RoPE}, no
biases) so it integrates into the open-source serving stack
unchanged \citep[\S 2]{ma2024-bitnet}. At 3B \glsterm{parameters}{parameters},
BitNet b1.58 matches FP16 LLaMA \glsterm{perplexity}{perplexity} ($9.91$ vs.\ $10.04$) at
$2.71\times$ lower latency and $3.55\times$ lower memory
\citep[\S 3, Table~1]{ma2024-bitnet}.\footnote{KB excerpt:
\texttt{kb/excerpts/ma2024-bitnet.md\#sec-3-table-1}.}

\begin{intuition}
Each regime operates on a different part of the inference graph.
SmoothQuant rescales activations and weights pre-runtime so a
hardware-friendly per-tensor INT8 GEMM suffices --- the MM kernel is
unchanged, the inputs are reshaped by \cref{eq:smoothquant-eq}. \glsterm{gptq}{GPTQ}
adjusts only the weight matrix via second-order layer-wise
reconstruction in \cref{eq:gptq-objective} --- activations stay FP16,
the LM-head and attention paths are untouched. BitNet replaces
multiplications with sign-flip-plus-select inside the \glsterm{bitlinear}{BitLinear} layer
of \cref{eq:bitnet-quant} --- the architecture changes, the energy
cost on a 7nm chip drops $71\times$ on the matmul-arithmetic
component \citep[\S 3]{ma2024-bitnet}. Each regime answers a different
question about which part of the deployment graph is precious.
\end{intuition}

\begin{contradiction}
\glsterm{nvfp4}{NVFP4}'s production status at frontier scale is contested. NVIDIA's
Blackwell-launch materials and vendor blogs report \glsterm{nvfp4}{NVFP4} as
production-ready for inference of large models with quality matching
\glsterm{fp8}{FP8} / FP16 baselines; third-party reproduction as of 2026-Q2 is sparse,
with most published benchmarks coming from NVIDIA-collaborator labs.
By contrast, \glsterm{fp8}{FP8} production training has independent confirmation
through \glsterm{precision-pinning}{DeepSeek-V3}'s open release \citep[\S 2.5]{deepseek-v3} and
multiple downstream replications. The KB reflects this asymmetry: \glsterm{fp8}{FP8}
is a tier-A primary claim, \glsterm{nvfp4}{NVFP4} is a tier-B forum-signal claim
pending peer-reviewed third-party benchmarks.
\deepencite{NVFP4 frontier-scale quality}
\end{contradiction}

\subsection{Structured output as constrained decoding}
\label{sec:constrained-decoding}

The third inference-time concern is enforcing that generated text
conforms to a target language $\mathcal{G}$ (regex, JSON schema, \glsterm{gbnf}{GBNF},
arbitrary CFG). The mechanism is \emph{logit masking}: at each step,
zero out the probability of any token whose acceptance would leave the
partial output outside $\mathcal{G}$, then sample from the renormalized
distribution. Formally, if $\ell_t \in \R^V$ are the model \glsterm{logits}{logits} at
step $t$ and $s_t$ is the grammar-state (FSM state, parser stack
summary), the constrained-decoding distribution is
\begin{equation}
P(y_t = v \mid y_{<t}) = \frac{\mathbb{1}[v \in \mathcal{V}_{\text{allowed}}(s_t)] \cdot \exp(\ell_{t,v})}{\sum_{v' \in \mathcal{V}_{\text{allowed}}(s_t)} \exp(\ell_{t,v'})},
\label{eq:constrained-decoding}
\end{equation}
implemented by setting $\ell_{t,v} \leftarrow -\infty$ for
$v \notin \mathcal{V}_{\text{allowed}}(s_t)$ before sampling. The
expensive operation is computing
$\mathcal{V}_{\text{allowed}}(s_t) \subseteq \{1, \dots, V\}$ at every
step: a naive grammar engine queries acceptance for every token in a
$100$K-vocab model, costing $\sim 100$ms/step, which dominates the
forward pass.

\citet{willard2023-outlines} reformulate the problem as
\emph{transitions in an FSM whose alphabet is the model's \glsterm{vocabulary}{vocabulary}}.
The standard FSM operates over bytes; the contribution is a
precomputed \emph{index} mapping each FSM state $s$ to the set of
\glsterm{vocabulary}{vocabulary} tokens whose byte sequence keeps the FSM in a valid path
\citep[\S 3]{willard2023-outlines}\footnote{KB excerpt:
\texttt{kb/excerpts/willard2023-outlines.md\#sec-3}.}:
\begin{equation}
\mathcal{V}_{\text{allowed}}(s) = \{v \in \mathcal{V} : \mathrm{FSM} \text{ accepts the byte sequence of } v \text{ starting from } s\}.
\label{eq:outlines-index}
\end{equation}
The index is built once per grammar; per-step lookup is $O(1)$
amortized. Extension to context-free grammars goes via LALR(1) parsers
where the parser state and a stack-summary jointly \glsterm{key}{key} the index
\citep[\S 4]{willard2023-outlines}; for JSON, Python, and SQL the
stack-summary stays small enough to keep the index practical.

\citet{xgrammar2024} extend Willard--Louf with a partition that
exploits a property of typical schemas: most \glsterm{vocabulary}{vocabulary} tokens'
allowed/disallowed status depends only on the current pushdown-automaton
state, not on stack contents \citep[\S 3]{xgrammar2024}.\footnote{KB
excerpt: \texttt{kb/excerpts/xgrammar2024.md\#sec-3}.} \glsterm{context-independent-context-dependent-vocabulary-split}{XGrammar}
precomputes a per-state \emph{bitmask} over $\mathcal{V}$ for the
context-independent fraction (typically $\sim 99\%$ on JSON schemas);
the small context-dependent set ($\sim 1\%$, e.g.\ closing brackets
that must match the correct opening) is checked at runtime against a
persistent stack \citep[\S 4--5]{xgrammar2024}. The two-class engine
runs grammar checks on the CPU concurrently with the GPU forward pass
\citep[\S 6]{xgrammar2024}, achieving up to $100\times$ speedup over
prior solutions and $\leq 40\,\mu$s per token of grammar overhead on
$100$K-vocab models \citep[\S 7]{xgrammar2024}. By 2026 \glsterm{context-independent-context-dependent-vocabulary-split}{XGrammar} is
the default constrained-decoding backend in vLLM, \glsterm{compressed-fsm-2}{SGLang}, and
TensorRT-\glsterm{llm}{LLM} \citep[FORUM-SIGNAL]{xgrammar2024}.

\begin{analogy}
\glsterm{constrained-decoding-2}{Constrained decoding} is the \emph{product} of the FSM (which says
which tokens are allowed) with the language model (which says which
tokens are likely): the masked distribution in
\cref{eq:constrained-decoding} is the model's \glsterm{softmax}{softmax} times the
indicator function of $\mathcal{V}_{\text{allowed}}(s_t)$,
renormalized. This is why it is architectural rather than a sampling
choice: the LM head's columns must align one-to-one with the tokenizer
that the FSM compiles against \cref{eq:outlines-index}, and the head's
output dimension $V$ is the same $V$ over which the bitmask is stored.
A model whose tokenizer or LM-head is replaced post-hoc invalidates
every precomputed index. The choice in
\cref{sec:tokenization-embedding} forward-determines what the
constrained-decoding stack can do.
\end{analogy}

\subsection{Forward link}

The three concerns in this section --- KV-block layout, quantization
regime, constrained-decoding stack --- jointly with the architectural
choices of \cref{sec:tokenization-embedding} through
\cref{sec:multimodal} fix the deployable surface of a 2026 frontier
\glsterm{llm}{LLM}. We turn next (\cref{sec:frontier-snapshot}) to a model-by-model
table of the architectural choices the leading 2026 systems share, and
where they diverge. \glsterm{pagedattention}{PagedAttention} has become universal; \glsterm{gqa}{GQA}-or-\glsterm{multi-head-latent-attention}{MLA} is
near-universal; ternary pretraining at 70B+ remains open; \glsterm{nvfp4}{NVFP4}
production at frontier scale remains contested. The localized open
frontiers we have surfaced through \cref{sec:kv-sharing,sec:hardware-attention,sec:state-space}
and the ones surfaced here are the ones \cref{sec:contradictions} will
address as a single concentrated discussion.
